% =============================================================
%  IEEE Conference paper
%  Compile with:  pdflatex full_paper.tex   (twice for refs)
% =============================================================
\documentclass[conference,a4paper]{IEEEtran}
\IEEEoverridecommandlockouts

% --- packages ---
\usepackage{cite}
\usepackage{amsmath,amssymb,amsfonts}
\usepackage{algorithmic}
\usepackage{graphicx}
\usepackage{textcomp}
\usepackage{xcolor}
\usepackage{booktabs}
\usepackage{array}
\usepackage{url}
\usepackage[hidelinks]{hyperref}

\def\BibTeX{{\rm B\kern-.05em{\sc i\kern-.025em b}\kern-.08em
    T\kern-.1667em\lower.7ex\hbox{E}\kern-.125emX}}

\begin{document}

% =============================================================
%  TITLE  (pick one; delete others before submission)
% =============================================================
\title{A Calibrated Stacked-Ensemble Approach to IMDb Rating
       Prediction from Movie Screenplays}
% Alternate title 1: How Much Does the Screenplay Actually Predict?
%                    Calibrated Baselines for IMDb Rating Prediction
% Alternate title 2: Stacked SBERT--XGBoost for Screenplay-Based
%                    IMDb Rating Prediction with Honest Statistical Reporting

% =============================================================
%  AUTHORS  (fill in)
% =============================================================
\author{
\IEEEauthorblockN{1\textsuperscript{st} Given Name Surname}
\IEEEauthorblockA{\textit{dept.\ name of organization} \\
\textit{name of organization (of Affiliation)}\\
City, Country \\
email@address \;or\; ORCID}
\and
\IEEEauthorblockN{2\textsuperscript{nd} Given Name Surname}
\IEEEauthorblockA{\textit{dept.\ name of organization} \\
\textit{name of organization (of Affiliation)}\\
City, Country \\
email@address \;or\; ORCID}
\and
\IEEEauthorblockN{3\textsuperscript{rd} Given Name Surname}
\IEEEauthorblockA{\textit{dept.\ name of organization} \\
\textit{name of organization (of Affiliation)}\\
City, Country \\
email@address \;or\; ORCID}
}

\maketitle

% =============================================================
%  ABSTRACT
% =============================================================
\begin{abstract}
We study how much of a movie's IMDb rating can be predicted from its
screenplay alone, using a corpus of 5{,}195 feature-film screenplays
joined with IMDb metadata. Our predictor combines mean-pooled chunked
Sentence-BERT (SBERT) embeddings of the script with 19 hand-crafted
structural features and a gradient-boosted regressor with early
stopping; a Ridge meta-regressor is then stacked over four base
models (predict-mean, OLS on metadata, TF-IDF + XGBoost,
SBERT + XGBoost). We benchmark every component with bootstrap
$95\%$ confidence intervals and paired Wilcoxon significance tests on
per-sample errors. Under five-fold cross-validation the stacked
system attains $\text{RMSE}=0.935\pm0.032$, $\text{MAE}=0.706\pm0.028$,
$R^{2}=0.577\pm0.010$, significantly improving over the strongest
single base model (paired Wilcoxon $p\approx 4\times 10^{-6}$) and
over every other baseline ($p\le 3\times 10^{-5}$). Two findings
deserve emphasis. First, three metadata features alone (year,
runtime, decade) explain $R^{2}\approx 0.38$, so the marginal
contribution of script content is $\Delta R^{2}\approx +0.18$ -
significant but smaller than the headline $R^{2}$ alone suggests.
Second, the legacy practice of inverse-frequency sample weighting
harms performance ($\Delta\text{MAE}=-0.048$,
$p\approx 3\times 10^{-26}$), contradicting earlier reports. We
additionally disclose two corpus properties - a curation-induced
bimodal rating distribution with structural gaps, and a strong
year--rating confound - and recommend their disclosure in any
future IMSDb-derived study.
\end{abstract}

\begin{IEEEkeywords}
IMDb rating prediction, screenplay analysis, Sentence-BERT,
stacked ensemble, gradient boosting, statistical reporting,
dataset bias.
\end{IEEEkeywords}

% =============================================================
%  I.  INTRODUCTION
% =============================================================
\section{Introduction}

Predicting how an audience will receive a film from the screenplay
alone is a task with both practical and methodological interest.
Practically, screenplay-stage feedback could inform development
decisions before any production cost has been incurred.
Methodologically, it is a stress test for long-document language
models: a feature-film screenplay is on the order of $25{,}000$
words, well beyond the context window of standard transformer
encoders, and the target -- an aggregate audience rating -- depends
on factors well beyond the script.

Three lines of prior work address this problem.
(i)~Classical text-feature pipelines treat the screenplay as a bag
of words or $n$-grams and feed the resulting sparse vectors to a
regressor \cite{ref3,ref4,ref26}; these approaches scale to long
documents but ignore semantics.
(ii)~Pre-trained sentence encoders such as SBERT \cite{ref2} produce
dense semantic embeddings; combined with chunking, they can
represent documents that exceed their native context window
\cite{ref19}.
(iii)~End-to-end long-document transformers such as
Longformer \cite{ref1} can process sequences up to a few thousand
tokens but require fine-tuning and GPU compute.

This paper focuses on~(ii) -- a frozen SBERT encoder combined with
engineered features and a gradient-boosted regressor -- and asks:
how does this hybrid compare to \emph{carefully chosen
non-transformer baselines}, what is the marginal contribution of
the SBERT component once metadata is accounted for, and which
training conventions in published prior work generalize to a
properly controlled comparison?

We organize the paper around three research questions:
\begin{itemize}
\item \textbf{RQ1.} Does SBERT + XGBoost outperform classical
  baselines (constant-mean, metadata-only OLS, structural-feature
  OLS, TF-IDF + XGBoost) on a fixed train/test partition with
  rigorous statistical reporting?
\item \textbf{RQ2.} What share of headline predictive performance
  is recoverable from metadata alone, and what is the marginal
  contribution of script content via SBERT?
\item \textbf{RQ3.} Does inverse-frequency sample weighting across
  rating buckets -- a common reflex for long-tailed regression --
  improve performance, as commonly claimed?
\end{itemize}

\textbf{Contributions.}
We make five contributions: (1) a rigorous baseline suite with
bootstrap CIs and paired Wilcoxon tests; (2) a stacked-ensemble
predictor that significantly improves over its strongest base model
with \emph{tighter} per-fold variance; (3) a decomposition of
headline $R^{2}$ between metadata and script content; (4) a negative
result on inverse-frequency sample weighting; and (5) an honest
characterization of an IMSDb-style screenplay corpus, including
structural rating gaps and a year--rating confound. A reproducible
pipeline (\texttt{experiments.py}, \texttt{stack.py},
\texttt{tune\_xgb.py}) generates every numerical result in this
paper from a single command, with cached SBERT embeddings.

% =============================================================
%  II.  RELATED WORK
% =============================================================
\section{Related Work}

\textbf{Movie rating and box-office prediction from text.}
A long line of work predicts box office or audience ratings from
script-derived features
\cite{ref3,ref4,ref7,ref8,ref9,ref15,ref17,ref26,ref27,ref29,ref33,ref37}.
Eliashberg et al.~\cite{ref3} applied kernel-based methods directly
to scripts; Hunter et al.~\cite{ref4} used document-frequency text
features; Kim et al.~\cite{ref9} predicted success from plot
summaries with deep models; Bristi et al.~\cite{ref26} surveyed
classical regressors; Cini~\cite{ref7} combined NLP with production
data; Pal et al.~\cite{ref37} focused on genre composition.
Joshi et al.~\cite{ref20} used pre-release critique text. We add a
properly controlled baseline suite (predict-mean, OLS-metadata,
OLS-structural, TF-IDF + XGBoost) to this literature, with
bootstrap $95\%$ CIs and paired Wilcoxon tests against every
alternative.

\textbf{Long-document NLP.}
Beltagy et al.'s Longformer~\cite{ref1} introduced sparse attention
for documents up to $4{,}096$ tokens; Chalkidis et al.~\cite{ref19}
explored hierarchical attention transformers for documents that
still exceed transformer context. Reimers and Gurevych's
Sentence-BERT~\cite{ref2} established that pretrained sentence
encoders can be combined with simple aggregation to represent
arbitrary-length text. Tixier~\cite{ref36} surveys the design
space. Our pipeline follows this hybrid path: chunked SBERT with
mean pooling.

\textbf{Affect, narrative, and character analysis from screenplays.}
Reagan et al.~\cite{ref5} characterized emotional arcs of stories;
Ramakrishna et al.~\cite{ref6} studied character portrayal
differences; Shafaei et al.~\cite{ref10} predicted MPAA ratings
from dialogue; Zhang et al.~\cite{ref11} graded severity in
screenplays; Chu et al.~\cite{ref12}, Hipson et al.~\cite{ref13},
and Elkins~\cite{ref14} studied emotion dynamics in dialogue and
narrative; Naeem et al.~\cite{ref18} applied sentiment analysis to
reviews; Kar et al.~\cite{ref35} tagged movies via plot synopsis
emotion flow. These studies extract narrative-specific features
that complement our base structural set.

\textbf{Movie data and external signals.}
Madongo et al.~\cite{ref16} mined trailers via RNNs;
Asur and Huberman~\cite{ref21}, Oghina et al.~\cite{ref22},
Mishne and Glance~\cite{ref23}, and Mest{\'y}an et al.~\cite{ref24}
used social media or Wikipedia activity; Balestri et al.~\cite{ref31}
generated trailers via LLMs; Sharma et al.~\cite{ref32} released a
larger movie dataset; Mohamed et al.~\cite{ref34} released an
age-appropriateness dataset. Vogel~\cite{ref25} provides industry
context. Our work uses screenplay text + IMDb metadata only;
integrating these external signals is a natural extension.

% =============================================================
%  III.  DATASET
% =============================================================
\section{Dataset}

\textbf{Source.}
We construct the corpus by joining (i)~full screenplay texts
collected from the Internet Movie Script Database (IMSDb) and
(ii)~movie metadata (release year, runtime, IMDb rating) retrieved
from IMDb. The joined dataset consists of $5{,}204$ records, of
which $5{,}195$ have script files of at least $1$ KB and a valid
IMDb rating after cleaning.

\textbf{Per-record fields.}
Movie name, year ($1922$--$2025$ after correcting four typographic
errors), decade (label-encoded), IMDb rating ($1.0$--$10.0$, one
decimal), IMDb ID, runtime ($45$--$254$ min, median $99$), and the
screenplay file (median $44.7$ KB; 5th percentile $19$ KB).

\begin{table}[!t]
\caption{Bucketed Rating Counts ($n = 5{,}195$).}
\label{tab:buckets}
\centering
\begin{tabular}{@{}lrr@{}}
\toprule
Range & $n$ & \% \\
\midrule
Low $[1, 4)$        & $555$  & $10.7$ \\
Medium $[4, 6)$     & $2{,}736$ & $52.6$ \\
Good $[6, 8)$       & $1{,}685$ & $32.4$ \\
Excellent $[8, 10)$ & $228$  & $4.4$ \\
\bottomrule
\end{tabular}
\end{table}

Mean $5.98$, median $5.80$, std $1.44$, IQR $5.30$--$7.40$.

\textbf{Sampling artifact (important caveat).}
Enumerating every unique rating value in the dataset reveals
\emph{structural gaps}: zero records in $[4.3, 5.0)$ and zero
records in $[6.0, 7.0)$. Together these intervals cover
approximately $17\%$ of typical IMDb rating mass. We attribute this
to selection effects upstream of our control: the IMSDb collection
is editorially curated, producing a corpus that is
\emph{bimodal by construction}. All metrics in this paper should be
interpreted as conditional on IMSDb-style films, not arbitrary
IMDb films.

\textbf{Temporal coverage and confound.}
Mean rating decreases monotonically with release decade (1930s
mean $7.49$ $\rightarrow$ 2010s mean $5.48$; Spearman
$\rho \approx -0.93$ over decade midpoints). We attribute this to
survivorship bias -- only canonical older films appear in IMSDb --
combined with the modern sample being large enough to cover the
full quality spectrum. Any model with access to \texttt{Year} or
\texttt{Decade} features can exploit this calendar confound; we
therefore present a metadata-only OLS baseline in
Section~\ref{sec:results} to quantify the share of rating variance
that this confound alone explains.

% =============================================================
%  IV.  METHODOLOGY
% =============================================================
\section{Methodology}

\textbf{Preprocessing.}
For each raw screenplay we produce two text variants: an
aggressive \texttt{clean\_text} (lowercased; stage directions,
character cues, INT./EXT. markers, timestamps, scene numbers, and
most punctuation removed) used as input to TF-IDF and structural-
feature extraction; and a light \texttt{clean\_text\_for\_sbert}
that preserves case, punctuation, and paragraph structure -- these
properties are part of SBERT's pretraining surface and are degraded
by aggressive normalization.

\textbf{Hand-crafted structural features ($n = 19$).}
Length / volume (\texttt{char\_count}, \texttt{word\_count},
\texttt{line\_count}, \texttt{sentence\_count}); vocabulary
complexity (\texttt{avg\_word\_length}, \texttt{unique\_word\_ratio},
\texttt{long\_word\_ratio}); sentence structure
(\texttt{avg\_sentence\_length}, \texttt{sentence\_length\_std});
dialogue (\texttt{dialogue\_density}, \texttt{unique\_characters});
emotional indicators (\texttt{exclamation\_ratio},
\texttt{question\_ratio}); action / structure
(\texttt{action\_density}, \texttt{scene\_count},
\texttt{words\_per\_scene}); metadata (\texttt{year},
\texttt{decade\_encoded}, \texttt{movie\_length}). Missing values
are imputed with the training-set median; features are standardized
to zero mean and unit variance via a \texttt{StandardScaler} fit on
the training fold only.

\textbf{SBERT encoding with chunking.}
We use \texttt{all-MiniLM-L6-v2}~\cite{ref2}, which produces
$384$-dim document embeddings. To handle screenplays exceeding
SBERT's $256$-token context window, each document is chunked into
$256$-word windows with a $50$-word overlap; per-chunk embeddings
are mean-pooled to a single $384$-dim document embedding.
Per-document embeddings are computed once for the corpus and cached
on disk, keyed by \texttt{(model\_name, n\_scripts, chunk\_size,
overlap, pooling, preprocessing\_version)}; subsequent experiments
reuse the cache.

\textbf{Main system.}
The main predictor concatenates the $384$-dim SBERT embedding with
the standardized $19$-dim feature vector, producing a $403$-dim
input fed to an XGBoost regressor (\texttt{max\_depth=6},
\texttt{learning\_rate=0.05}, \texttt{reg\_alpha=0.1},
\texttt{reg\_lambda=1.0}, \texttt{tree\_method=hist},
\texttt{n\_estimators=1000} capped by
\texttt{early\_stopping\_rounds=20} on a held-out validation
split). We do \emph{not} use inverse-frequency sample weights; the
ablation in Section~\ref{sec:results} shows they degrade
performance.

\textbf{Baselines.}
(1)~\texttt{predict\_mean}: constant-mean predictor -- the floor.
(2)~\texttt{ols\_metadata}: OLS on
$[\texttt{year}, \texttt{decade\_encoded}, \texttt{movie\_length}]$
only.
(3)~\texttt{ols\_structural}: OLS on the full $19$-dim structural
feature vector.
(4)~\texttt{tfidf\_xgboost}: TF-IDF (top $8{,}000$ uni- and
bigrams, \texttt{min\_df=3}, \texttt{max\_df=0.85}, sublinear TF) +
XGBoost. All baselines train and evaluate on the same
train/val/test splits as the main system; predictions are clipped
to $[1, 10]$.

\textbf{Stacked ensemble.}
A Ridge meta-regressor (\texttt{alpha=1.0}) is fit over the four
base models' predictions using nested cross-validation: an outer
$5$-fold KFold partitions the corpus into disjoint test folds;
within each outer training portion, an inner $5$-fold KFold
produces \emph{out-of-fold} base-model predictions on which the
meta-regressor is trained; the base models are then refit on the
full outer training set and evaluated on the outer test fold via
the trained meta-regressor.

\textbf{Hyperparameter search.}
We additionally run a $100$-trial Optuna TPE search over the
XGBoost head's hyperparameter space (learning rate, depth, min
child weight, subsample, column-subsample, $L_{1}$/$L_{2}$
regularization, gamma) with \texttt{early\_stopping\_rounds=30} on
a fixed train/val split internal to the training partition.

\textbf{Splits and statistical protocol.}
Headline numbers use a $70$/$15$/$15$ train/validation/test
partition with \texttt{random\_state=42}. Robustness numbers use
$5$-fold cross-validation. We report bootstrap $95\%$ confidence
intervals on RMSE/MAE/$R^{2}$ ($1{,}000$ resamples). For
comparisons against the main system on the same test samples, we
additionally report a paired-bootstrap $95\%$ CI for
$\Delta\text{MAE}$ and a paired Wilcoxon signed-rank test on
per-sample absolute errors.

% =============================================================
%  V.  RESULTS
% =============================================================
\section{Results}
\label{sec:results}

\begin{table*}[!t]
\caption{Single-Split Test Metrics with $95\%$ Bootstrap CIs ($n_{\text{test}} = 780$).}
\label{tab:single-split}
\centering
\begin{tabular}{@{}lccc@{}}
\toprule
Model & RMSE & MAE & $R^{2}$ \\
\midrule
\texttt{predict\_mean}             & $1.518$ $[1.46, 1.58]$ & $1.247$ $[1.19, 1.31]$ & $-0.001$ \\
\texttt{ols\_metadata}             & $1.184$ $[1.13, 1.24]$ & $0.946$ $[0.89, 0.99]$ & $0.391$ $[0.35, 0.43]$ \\
\texttt{ols\_structural}           & $1.129$ $[1.07, 1.19]$ & $0.881$ $[0.83, 0.93]$ & $0.447$ $[0.41, 0.49]$ \\
\texttt{tfidf\_xgboost}            & $1.129$ $[1.06, 1.20]$ & $0.854$ $[0.80, 0.90]$ & $0.447$ $[0.39, 0.50]$ \\
\texttt{sbert\_xgboost} (weighted) & $1.032$ $[0.97, 1.09]$ & $0.789$ $[0.74, 0.84]$ & $0.538$ $[0.48, 0.59]$ \\
\textbf{\texttt{sbert\_xgboost} (no weights)}
                                   & $\mathbf{0.997}$ $[0.94, 1.06]$
                                   & $\mathbf{0.749}$ $[0.70, 0.79]$
                                   & $\mathbf{0.568}$ $[0.53, 0.61]$ \\
\bottomrule
\end{tabular}
\end{table*}

\begin{table}[!t]
\caption{Paired Wilcoxon vs.\ SBERT (no weights), Single Split.
Positive $\Delta$MAE means SBERT wins.}
\label{tab:paired-single}
\centering
\begin{tabular}{@{}lrlc@{}}
\toprule
Baseline & $\Delta$MAE & 95\% CI & Wilcoxon $p$ \\
\midrule
\texttt{predict\_mean}   & $+0.458$ & $[+0.39, +0.52]$  & $4.3\!\times\!10^{-34}$ \\
\texttt{ols\_metadata}   & $+0.157$ & $[+0.11, +0.20]$  & $9.1\!\times\!10^{-11}$ \\
\texttt{ols\_structural} & $+0.093$ & $[+0.05, +0.13]$  & $2.7\!\times\!10^{-5}$ \\
\texttt{tfidf\_xgboost}  & $+0.065$ & $[+0.015, +0.11]$ & $1.4\!\times\!10^{-2}$ \\
\bottomrule
\end{tabular}
\end{table}

\textbf{Cross-validated robustness.}
Table~\ref{tab:cv} reports the $5$-fold CV ordering, which preserves
every conclusion above with tighter confidence; the previously
borderline win over TF-IDF resolves at $p \approx 2.6\times 10^{-5}$
under pooled paired Wilcoxon ($n \approx 5{,}195$).

\begin{table}[!t]
\caption{$5$-Fold CV Test Metrics (mean $\pm$ std).}
\label{tab:cv}
\centering
\begin{tabular}{@{}lccc@{}}
\toprule
Model & RMSE & MAE & $R^{2}$ \\
\midrule
\texttt{predict\_mean}             & $1.438\pm0.056$ & $1.175\pm0.057$ & $0.000\pm0.000$ \\
\texttt{ols\_metadata}             & $1.130\pm0.047$ & $0.883\pm0.042$ & $0.382\pm0.017$ \\
\texttt{ols\_structural}           & $1.076\pm0.040$ & $0.833\pm0.031$ & $0.440\pm0.007$ \\
\texttt{tfidf\_xgboost}            & $1.083\pm0.034$ & $0.819\pm0.022$ & $0.433\pm0.014$ \\
\texttt{sbert\_xgboost} (weighted) & $1.000\pm0.027$ & $0.768\pm0.024$ & $0.516\pm0.026$ \\
\texttt{sbert\_xgboost} (no weights) & $0.958\pm0.033$ & $0.720\pm0.027$ & $0.556\pm0.009$ \\
\bottomrule
\end{tabular}
\end{table}

\textbf{Sample-weighting ablation.}
In both regimes, removing inverse-frequency sample weights
\emph{significantly improves} performance: single-split
$\Delta\text{MAE} = -0.040$ ($[-0.07, -0.01]$), Wilcoxon
$p = 1.2\times 10^{-2}$; $5$-fold pooled
$\Delta\text{MAE} = -0.048$ ($[-0.06, -0.04]$), Wilcoxon
$p = 2.5\times 10^{-26}$. The $13.24\times$ weight on the small
``Excellent'' bucket destabilizes XGBoost; we adopt the unweighted
variant as the headline configuration.

\textbf{Metadata decomposition.}
\texttt{ols\_metadata} alone -- three numbers -- attains
$R^{2} = 0.391$ (single split) and $0.382 \pm 0.017$
($5$-fold). Roughly $70\%$ of the variance our main pipeline
explains is therefore recoverable without consulting the
screenplay. The marginal contribution of script content (SBERT +
structural features over metadata-only) is
$\Delta R^{2} \approx +0.18$, statistically robust under pooled
Wilcoxon ($p \approx 7\times 10^{-37}$) but smaller than the
headline figure suggests.

\textbf{Stacked ensemble.}
A Ridge meta-regressor over the four base models, trained via
nested CV (Section~IV), lifts performance further
(Table~\ref{tab:stack}).

\begin{table}[!t]
\caption{$5$-Fold CV Test Metrics: Stacked Ensemble vs.\ Strongest
Base Model (Same Outer Folds as Table~\ref{tab:cv}).}
\label{tab:stack}
\centering
\begin{tabular}{@{}lccc@{}}
\toprule
Model & RMSE & MAE & $R^{2}$ \\
\midrule
\texttt{sbert\_xgboost} (no weights) & $0.956\pm0.034$ & $0.719\pm0.028$ & $0.558\pm0.013$ \\
\textbf{stacked (Ridge over 4)}      & $\mathbf{0.935\pm0.032}$ & $\mathbf{0.706\pm0.028}$ & $\mathbf{0.577\pm0.010}$ \\
\bottomrule
\end{tabular}
\end{table}

Pooled paired Wilcoxon stacked vs.\ SBERT:
$\Delta\text{MAE} = +0.013$ ($[+0.008, +0.019]$),
$p = 4.2\times 10^{-6}$. The stacked $R^{2}$ standard deviation
($0.010$) is \emph{tighter} than any single base model ($0.013$
for SBERT, $0.014$ for TF-IDF). Across all five folds the Ridge
meta-regressor's coefficients are remarkably stable -- SBERT
$\approx 0.69$, TF-IDF $\approx 0.38$, \texttt{ols\_metadata}
$\approx 0.27$, \texttt{ols\_structural} $\approx -0.13$. The
negative coefficient on \texttt{ols\_structural} indicates its
information is fully absorbed by upstream models once SBERT and
TF-IDF predictions are available.

\textbf{Hyperparameter tuning (Optuna).}
A $100$-trial TPE search on the XGBoost head improves single-split
test $R^{2}$ from $0.568$ to $0.581$ $[0.55, 0.62]$
($\Delta\text{MAE} = -0.035$). The best configuration uses
\texttt{learning\_rate} $\approx 0.011$ ($5\times$ lower than the
hand-picked default), \texttt{reg\_alpha} $\approx 1.0$
($10\times$ higher), \texttt{max\_depth}$=5$ -- directly addressing
the overfitting visible in the legacy training curve.

\textbf{Per-bucket performance.}
Test-set MAE by rating bucket (no-weight SBERT):
Low $[1, 4)$ MAE $= 1.577$ ($n = 107$);
Medium $[4, 6)$ $0.558$ ($n = 380$);
Good $[6, 8)$ $0.672$ ($n = 250$);
Excellent $[8, 10)$ $0.818$ ($n = 43$). Performance is best on the
corpus mode and degrades on both tails -- predictions cluster
between approximately $4$ and $8$ even when true ratings span
$1.5$ to $9.3$.

% =============================================================
%  VI.  DISCUSSION
% =============================================================
\section{Discussion}

The headline $5$-fold $R^{2} \approx 0.58$ is consistent with a
moderately useful predictor of audience reception from screenplay
content. It is not, however, evidence that an audience rating is
``in the script'': three metadata features alone already explain
$R^{2} \approx 0.38$ on the same corpus and splits. The
$\Delta R^{2} \approx +0.18$ attributable to script content (SBERT
+ structural features) is statistically robust but places a hard
ceiling on how much an audience rating can be claimed to
``recover from the screenplay.''

Two corollaries follow. First, papers in this area should report
metadata-only OLS as a mandatory baseline; without it, reported
$R^{2}$ values systematically over-attribute predictive power to
the text component. Second, for downstream practitioners, the
\emph{delta} over a metadata baseline, not the absolute $R^{2}$, is
the relevant figure of merit.

The negative result on inverse-frequency sample weighting is also
notable. The most up-weighted bucket (Excellent, $13.24\times$
weight) contains only $147$ training samples; amplifying its
gradients makes the gradient-boosted regressor sensitive to a
high-leverage minority. Removing the reweighting also improves
per-bucket MAE on Low, suggesting the legacy weighting did not
deliver the rebalancing it was designed to produce.

% =============================================================
%  VII.  LIMITATIONS
% =============================================================
\section{Limitations}

\textbf{Selection bias.}
The corpus is editorially curated; rating intervals $[4.3, 5.0)$
and $[6.0, 7.0)$ are entirely empty. Generalization claims are
scoped to IMSDb-style films, not arbitrary IMDb draws.

\textbf{Metadata dominance.}
Our pipeline does not include a metadata-removed ablation isolating
the SBERT contribution from year/length/decade alone. The closest
available comparison (\texttt{ols\_structural}, which still
contains those three features) achieves $R^{2} = 0.44$; a
controlled metadata-removed ablation is the most important missing
experiment.

\textbf{Single encoder, single pooling.}
We report \texttt{all-MiniLM-L6-v2} with mean-pooling. A larger
encoder (\texttt{all-mpnet-base-v2}) and alternative pooling
operators (max, $\ell_{2}$-weighted) are queued; our refactored
embedding cache supports them without re-encoding.

\textbf{No properly-controlled long-document transformer baseline.}
A fair Longformer comparison (matched feature budget, full
$4{,}096$-token context) requires GPU resources beyond the present
scope.

\textbf{Rating subjectivity.}
IMDb ratings reflect production values, marketing, cast / director
reputation, and voter selection beyond the screenplay. The
unexplained variance ($1 - R^{2} \approx 0.42$) plausibly contains
a sizeable irreducible component.

% =============================================================
%  VIII.  CONCLUSION
% =============================================================
\section{Conclusion}

We presented a screenplay-to-IMDb-rating predictor combining frozen
SBERT embeddings, hand-crafted structural features, gradient
boosting, and a Ridge meta-regressor stacked over four base models,
and evaluated it under both single-split and $5$-fold CV protocols
with bootstrap CIs and paired Wilcoxon tests. The stacked system
attains $R^{2} = 0.577 \pm 0.010$ ($5$-fold CV), significantly
improving on the strongest single base model
($p \approx 4\times 10^{-6}$). We documented (i)~that approximately
$70\%$ of explained variance is recoverable from three metadata
features, (ii)~a negative result on inverse-frequency sample
weighting, (iii)~complementary signal across base models confirmed
by stacked variance, and (iv)~corpus selection artifacts that
warrant disclosure. The pipeline is reproducible from a single
command, runs in minutes on CPU, and produces a deployable $2$ MB
model.

Future work: a metadata-removed ablation; a properly-controlled
long-document transformer baseline; alternative pooling operators
using the released chunk-cache; evaluation on out-of-corpus
screenplays.

% =============================================================
%  REFERENCES
% =============================================================
\begin{thebibliography}{99}

\bibitem{ref1}
I.~Beltagy \emph{et al.}, ``Longformer: The long-document
transformer,'' \emph{arXiv preprint} arXiv:2004.05150, 2020.

\bibitem{ref2}
N.~Reimers and I.~Gurevych, ``Sentence-BERT: Sentence embeddings
using Siamese BERT-networks,'' in \emph{Proc.\ EMNLP}, 2019.

\bibitem{ref3}
J.~Eliashberg \emph{et al.}, ``Assessing box office performance
using movie scripts: A kernel-based approach,''
\emph{IEEE Trans.\ Knowl.\ Data Eng.}, 2014.

\bibitem{ref4}
S.~D.~Hunter \emph{et al.}, ``Predicting box office from the
screenplay: A text analytical approach,'' West East Institute,
2016.

\bibitem{ref5}
A.~J.~Reagan \emph{et al.}, ``The emotional arcs of stories are
dominated by six basic shapes,'' \emph{EPJ Data Science},
vol.~5, no.~31, 2016.

\bibitem{ref6}
A.~Ramakrishna \emph{et al.}, ``Linguistic analysis of differences
in portrayal of movie characters,'' in \emph{Proc.\ ACL}, 2017.

\bibitem{ref7}
K.~Cini, ``Forecasting film audience ratings: A natural language
processing approach to script and production data,''
\emph{Entertainment Computing}, 2025.

\bibitem{ref8}
J.~A.~Gross and T.~Roberson, ``Film success prediction using NLP
techniques,'' Stanford CS230 Project, 2021.

\bibitem{ref9}
Y.~J.~Kim \emph{et al.}, ``Prediction of movie success from plot
summaries using deep learning,'' in \emph{ACL Workshop}, 2019.

\bibitem{ref10}
M.~Shafaei \emph{et al.}, ``Age suitability rating: Predicting
MPAA rating based on movie dialogues,'' in \emph{LREC}, 2020.

\bibitem{ref11}
Y.~Zhang \emph{et al.}, ``From none to severe: Predicting severity
in movie scripts,'' in \emph{Findings of EMNLP}, 2021.

\bibitem{ref12}
E.~Chu, D.~Roy, and J.~Glass, ``Audio-visual sentiment analysis
for learning emotional arcs in movies,'' in \emph{ICCV}, 2017.

\bibitem{ref13}
W.~E.~Hipson \emph{et al.}, ``Emotion dynamics in movie
dialogues,'' \emph{PLoS ONE}, 2021.

\bibitem{ref14}
K.~Elkins, ``Beyond plot: How sentiment analysis reshapes
narrative structure,'' \emph{Cultural Analytics}, 2025.

\bibitem{ref15}
R.~Sharda and D.~Delen, ``Predicting box-office success of motion
pictures with neural networks,'' \emph{Expert Systems with
Applications}, 2006.

\bibitem{ref16}
C.~T.~Madongo \emph{et al.}, ``Movie box-office revenue prediction
model by mining deep features from trailers using recurrent neural
networks,'' \emph{J.\ of Advances in Information Technology},
2024.

\bibitem{ref17}
A.~Bhadrashetty and S.~Patil, ``Movie success and rating
prediction using data mining,'' \emph{J.\ of Scientific Research
and Technology}, 2024.

\bibitem{ref18}
M.~Z.~Naeem \emph{et al.}, ``Classification of movie reviews using
sentiment analysis,'' 2022.

\bibitem{ref19}
I.~Chalkidis \emph{et al.}, ``An exploration of hierarchical
attention transformers,'' \emph{arXiv preprint}, 2022.

\bibitem{ref20}
D.~Joshi \emph{et al.}, ``Pre-release critique text features for
revenue prediction.''

\bibitem{ref21}
S.~Asur and B.~A.~Huberman, ``Predicting the future with social
media,'' in \emph{WWW}, 2010.

\bibitem{ref22}
A.~Oghina \emph{et al.}, ``Predicting IMDb movie ratings using
Twitter,'' 2012.

\bibitem{ref23}
G.~Mishne and N.~Glance, ``Predicting movie sales from blogger
sentiment,'' in \emph{AAAI Spring Symposium}, 2006.

\bibitem{ref24}
M.~Mest{\'y}an \emph{et al.}, ``Early prediction of movie box
office success based on Wikipedia activity big data,''
\emph{PLoS ONE}, 2013.

\bibitem{ref25}
H.~L.~Vogel, \emph{Entertainment Industry Economics}, 10th ed.
Cambridge University Press, 2020.

\bibitem{ref26}
W.~R.~Bristi \emph{et al.}, ``Predicting IMDb rating of movies by
machine learning techniques,'' in \emph{IEEE}, 2019.

\bibitem{ref27}
A.~L.~Gomes \emph{et al.}, ``Predicting IMDb rating of TV series
with deep learning: The case of Arrow,''
\emph{arXiv preprint}, 2022.

\bibitem{ref28}
J.~Ramos \emph{et al.}, ``Movie rating prediction using sentiment
features,'' in \emph{SALLD}, 2022.

\bibitem{ref29}
V.~Udandarao \emph{et al.}, ``Movie revenue prediction using
machine learning models,'' \emph{arXiv preprint}, 2024.

\bibitem{ref30}
W.~Xie \emph{et al.}, ``Predicting movie success with multi-task
learning: GPT-based sentiment and SIR propagation,''
\emph{arXiv preprint}, 2025.

\bibitem{ref31}
R.~Balestri \emph{et al.}, ``An automatic deep learning approach
for trailer generation through large language models,''
\emph{arXiv preprint}, 2026.

\bibitem{ref32}
A.~S.~Sharma \emph{et al.}, ``Presenting a larger up-to-date movie
dataset,'' \emph{arXiv preprint}, 2021.

\bibitem{ref33}
Y.~Ding \emph{et al.}, ``A machine learning model based on
data-driven movie derivatives market prediction,''
\emph{arXiv preprint}, 2022.

\bibitem{ref34}
E.~Mohamed \emph{et al.}, ``A first dataset for film age
appropriateness investigation,'' in \emph{LREC}, 2020.

\bibitem{ref35}
S.~Kar \emph{et al.}, ``Folksonomication: Predicting tags for
movies from plot synopses using emotion flow encoded neural
network,'' in \emph{COLING}, 2018.

\bibitem{ref36}
A.~J.~P.~Tixier, ``Notes on deep learning for NLP,''
\emph{arXiv preprint}, 2018.

\bibitem{ref37}
A.~Pal \emph{et al.}, ``Identifying movie genre compositions using
neural networks,'' in \emph{IEEE}, 2020.

\bibitem{ref38}
M.~C.~Chiu \emph{et al.}, ``Screenplay quality assessment: Can we
predict who gets nominated?'' NUSe, 2020.

\end{thebibliography}

\end{document}
